\begin{center}
\hypertarget{native-056-article038}{}\pdfbookmark[2]{38.}{article-038}38.
\end{center}
Haud inutile erit, ea quae in art. praec. sunt tradita exemplo illustrare. Proposita sit anomalia vera $=100^0$, excentricitas $=0{,}967\,6456\,7$, $\log q=9{,}765\,6500$. Ecce iam calculum pro $E$, $B$, $A$ et $t$:
\begin{center}
\begin{tabular}{@{}ll@{}}
\makebox[18em][l]{$\log\operatorname{tang}\tfrac{1}{2}v$\hfill .\,.\,.\,.} & $0{,}076\,1865$\\[.3em]
\makebox[18em][l]{$\log\BareRoot{\dfrac{1-e}{1+e}}$\hfill .\,.\,.\,.\,.} & $9{,}107\,9927$\\[.4em]
\hline
\makebox[18em][l]{$\log\operatorname{tang}\tfrac{1}{2}E$\hfill .\,.\,.} & $9{,}184\,1792$,
\end{tabular}
\end{center}
\noindent unde $\tfrac{1}{2}E=8^0\,41^{\prime}\,19^{\prime\prime}{,}32$, atque $E=17^0\,22^{\prime}\,38^{\prime\prime}{,}64$. Huic valori ipsius $E$ respondet $\log B=0{,}000\,0040$; porro invenitur in partibus radii $E=0{,}303\,2928$, $\sin E=0{,}298\,6643$, unde $\tfrac{9}{20}E+\tfrac{1}{20}\sin E=0{,}151\,4150$, cuius logarithmus $=9{,}180\,1689$, adeoque $\log A^{\frac{1}{2}}=9{,}180\,1649$. Hinc deducitur per formulam 1) art. praec.
\begin{center}
\begin{tabular}{@{}lrl l@{}}
\multicolumn{3}{@{}l}{\makebox[22em][l]{$\log\dfrac{2Bq^{\frac{3}{2}}}{k\BareRoot{(1-e)}}$\hfill .\,.\,.\,.\,.\,.\,.\,.\,.}} & $2{,}458\,9614$\\[.4em]
\multicolumn{3}{@{}l}{\makebox[22em][l]{$\log A^{\frac{1}{2}}$\hfill .\,.\,.\,.\,.\,.\,.\,.\,.\,.\,.\,.}} & $9{,}180\,1649$\\[.3em]
\hline
$\log$ & $43{,}56386$ & $=$ & $1{,}639\,1263$\\[.3em]
 & $19{,}98014$ & & \\
\cline{2-2}
 & $63{,}54400$ & $=t$ &
\end{tabular}
\end{center}
\begin{center}
\begin{tabular}{@{}ll@{}}
\makebox[22em][l]{$\log\dfrac{2B(1+9e)}{15k}\left(\dfrac{q}{1-e}\right)^{\frac{3}{2}}$\hfill .\,.\,.} & $3{,}760\,1038$\\[.4em]
\makebox[22em][l]{$\log A^{\frac{3}{2}}$\hfill .\,.\,.\,.\,.\,.\,.\,.\,.\,.\,.\,.} & $7{,}540\,4947$\\[.3em]
\hline
$\log 19{,}98014$\hfill $=$ & $1{,}300\,5985$
\end{tabular}
\end{center}
\noindent Tractando idem exemplum secundum methodum vulgarem, invenitur $e\sin E$ in \SourceLeaf{057}secundis $=59610^{\prime\prime}{,}79=16^0\,33^{\prime}\,30^{\prime\prime}{,}79$, unde anomalia media $=49^{\prime}\,7^{\prime\prime}{,}85=2947^{\prime\prime}{,}85$. Hinc et ex $\log k\left(\dfrac{1-e}{q}\right)^{\frac{3}{2}}=1{,}666\,4302$ derivatur $t=63{,}54410$. Differentia, quae hic tantum est $10\,000^{\mathrm{ma}}$ pars unius diei, conspirantibus erroribus facile triplo vel quadruplo maior evadere potuisset.

Ceterum patet, solo adiumento talis tabulae pro $\log B$ etiam problema inversum omni praecisione solvi posse, determinando $E$ per tentamina repetita, ita ut valor ipsius $t$ inde calculatus cum proposito congruat. Sed haec operatio satis molesta foret: quamobrem iam ostendemus, quomodo tabulam auxiliarem multo commodius adornare, tentamina vaga omnino evitare, totumque calculum ad algorithmum maxime concinnum atque expeditum reducere liceat, qui nihil desiderandum relinquere videtur.
